\documentclass[conference]{IEEEtran}
\IEEEoverridecommandlockouts
% Package imports for IEEEtran compatibility
\usepackage{cite}
\usepackage{amsmath,amssymb,amsfonts}
\usepackage{algorithmic}
\usepackage{algorithm}
\usepackage{graphicx}
\usepackage{textcomp}
\usepackage{xcolor}
\usepackage{booktabs}
\usepackage{multirow}
\usepackage{tabularx}
\usepackage{url}
\usepackage{microtype}

\DeclareMathOperator{\Acc}{Acc}
\DeclareMathOperator{\AUC}{AUC}

\def\BibTeX{{\rm B\kern-.05em{\sc i\kern-.025em b}\kern-.08em
    T\kern-.1667em\lower.7ex\hbox{E}\kern-.125emX}}

\begin{document}

\title{ASUC-SOM: Subspace-Aware Adaptive Routing for Sequential Machine Unlearning under GDPR Article 17}

\author{\IEEEauthorblockN{Mithul Kannan}
\IEEEauthorblockA{\textit{dept. of CSE (Artificial Intelligence and Machine Learning)} \\
\textit{The National Institute Of Engineering}\\
Mysuru, Karnataka, India \\
mithulkannan17@gmail.com}
\and
\IEEEauthorblockN{Manveeth H V}
\IEEEauthorblockA{\textit{dept. of CSE (Artificial Intelligence and Machine Learning)} \\
\textit{The National Institute Of Engineering}\\
Mysuru, Karnataka, India \\
email address or ORCID}
\and
\IEEEauthorblockN{M M Kaverappa}
\IEEEauthorblockA{\textit{dept. of CSE (Artificial Intelligence and Machine Learning)} \\
\textit{The National Institute Of Engineering}\\
Mysuru, Karnataka, India \\
email address or ORCID}
\and
\IEEEauthorblockN{M Bharath Mani}
\IEEEauthorblockA{\textit{dept. of CSE (Artificial Intelligence and Machine Learning)} \\
\textit{The National Institute Of Engineering}\\
Mysuru, Karnataka, India \\
email address or ORCID}
\\[2.5ex]
\IEEEauthorblockN{Sumanth C S}
\IEEEauthorblockA{\textit{Assistant Professor} \\
\textit{dept. of CSE (Artificial Intelligence and Machine Learning)} \\
\textit{The National Institute Of Engineering}\\
Mysuru, Karnataka, India \\
email address or ORCID}
}

\maketitle

\begin{abstract}
Machine unlearning is predominantly evaluated under isolated, single-request assumptions: a model is trained, a single forget set is removed, and performance evaluation stops. Regulatory frameworks such as GDPR Article 17 and the California Consumer Privacy Act (CCPA) mandate continuous, sequential data erasure over the entire operational lifecycle of deployed systems. In this sequential regime, static unlearning heuristics (e.g., Selective Synaptic Dampening and Saliency Unlearning) suffer from severe parameter fatigue: repeated Fisher dampening indiscriminately degrades overlapping parameter dimensions, causing static SSD to catastrophically collapse to random guessing (9.96\% accuracy) by Round 3 under uniform random deletions. We present ASUC-SOM, an adaptive closed-loop controller that treats a stream of deletion requests as a geometric control problem. A Subspace Overlap Monitor (SOM) tracks an orthonormal basis $Q_k$ of historical gradient modification vectors, measuring projection overlap before weights are touched. An entropy-aware tiered router dynamically dispatches requests between rapid dampening (SSD, $\approx 150\times$ speedup), regularized saliency masking (SalUn, $\approx 115\times$ speedup), and selective exact retraining. A closed-loop verification gate enforces empirical forget and retain bounds with automated rollback and tier escalation. Evaluated on ResNet-18 across 20 sequential deletion rounds on CIFAR-10 (10,000 forgotten samples), ASUC-SOM maintains 97.19\% retain accuracy and 90.89\% generalization without experiencing rank collapse, delivering up to 148$\times$ compute speedup over full retraining.
\end{abstract}

\begin{IEEEkeywords}
machine unlearning, sequential deletion, GDPR Article 17, parameter fatigue, selective synaptic dampening, saliency, subspace monitoring, membership inference.
\end{IEEEkeywords}

\section{Introduction}
A trained deep neural network is, in a practical sense, an indexed compression of its training corpus. When an individual exercises the right to erasure under Article 17 of the European Union General Data Protection Regulation (GDPR) \cite{b1} or statutory deletion rights under the CCPA, the model operator must remove the influence of that data from the model weights. The cleanest answer is to retrain the model from scratch on the retained dataset without those records. It is also, for modern deep networks undergoing continuous deletion requests, an answer no operator can afford to execute repeatedly \cite{b2}, \cite{b3}. Machine unlearning grew out of this tension: the objective is to excise the influence of a forget set $D_f$ faster than full retraining, while proving that model amnesia was achieved.

The last few years have produced several effective answers to the single-request formulation of that question. Selective Synaptic Dampening (SSD) rescales parameters that are disproportionately sensitive to the forget set relative to the retain set, requiring no iterative gradient backpropagation \cite{b4}. Saliency Unlearning (SalUn) confines gradient ascent to parameters with top gradient saliency on the forget set \cite{b5}. Gradient ascent, Fisher-based noise injection \cite{b7}, and incompetent teacher distillation \cite{b10} fill out the approximate unlearning landscape. However, nearly all of these methods are evaluated identically: take one freshly trained model, remove one forget set, record post-unlearning metrics, and stop.

That evaluation protocol fails to capture how deletion occurs in production. Deletion requests arrive as a sequential stream over weeks and months, varying in batch cardinality and class distribution. The model processing the twentieth request is not the original checkpoint $\theta_0$; it is the cumulative product of nineteen prior parameter updates. 

Three structural failure modes arise when static unlearning methods are applied repeatedly:
\begin{enumerate}
    \item \textbf{Parameter Fatigue and Representational Collapse:} Static parameter updates accumulate directional distortion. On overlapping weight channels, repeated dampening causes catastrophic rank collapse. As demonstrated in our experiments, static SSD collapses to 9.96\% retain accuracy (random guessing on 10-class CIFAR-10) by Round 3.
    \item \textbf{Blind Heuristic Application:} No single unlearning primitive is optimal for all deletion geometries. Removing 50 balanced samples across all classes and removing 2,000 samples concentrated in a single class are geometrically distinct tasks, yet static algorithms treat them identically.
    \item \textbf{Absence of Closed-Loop Verification:} Static algorithms apply parameter updates without pre-checking subspace interference or post-verifying that amnesia and retain utility bounds were satisfied. Corrupted checkpoints are deployed into production undetected.
\end{enumerate}

To overcome these limitations, we introduce ASUC-SOM (Adaptive Sequential Unlearning Controller with Subspace Overlap Monitoring). Our contributions are:
\begin{itemize}
    \item \textbf{Subspace Overlap Monitor (SOM):} We formulate a geometric monitor that maintains an orthonormal basis $Q_k$ of historical update trajectories, computing gradient projection overlap before weights are modified to predict parameter fatigue.
    \item \textbf{Multi-Tier Dynamic Router:} We design a three-tier dispatch policy that maps SOM overlap scores and normalized class entropy $h$ to the optimal unlearning primitive (SSD, SalUn, or exact retraining).
    \item \textbf{Closed-Loop Self-Healing Gate:} We implement an automated verification gate with checkpoint rollback and tier escalation that prevents corrupted weights from leaving the system.
    \item \textbf{Comprehensive Sequential Benchmark:} Evaluated across 20 sequential deletion rounds on CIFAR-10 ResNet-18 (10,000 deleted samples), ASUC-SOM sustains 97.19\% retain accuracy and 90.89\% test accuracy without experiencing rank collapse.
\end{itemize}

\section{Related Work}

\subsection{Exact and Approximate Machine Unlearning}
Cao and Yang introduced the machine unlearning problem for statistical-query learners \cite{b2}, and Bourtoule et al. formulated SISA (Sharded, Isolated, Sliced, Aggregated) to make exact unlearning feasible by retraining isolated shards \cite{b3}. Ginart et al. developed deletion algorithms for $k$-means clustering \cite{b14}, and Sekhari et al. established generalization bounds for unlearning \cite{b15}. Approximate methods trade exact statistical equivalence for execution speed. Golatkar et al. scrub information using the Fisher Information Matrix \cite{b7}, Graves et al. undo specific SGD parameter updates \cite{b13}, and Chundawat et al. distil unlearning behavior from an incompetent teacher \cite{b10}. Foster et al. proposed SSD \cite{b4}, and Fan et al. introduced SalUn \cite{b5}. Thudi et al. showed that approximate unlearning metrics can often be brittle \cite{b8}.

\subsection{Sequential Deletion and Adaptive Requests}
Gupta et al. established theoretical bounds for adaptive unlearning where deletion requests may depend adversarially on earlier model outputs \cite{b12}. Our work focuses on the empirical systems engineering challenge: characterizing why fast unlearning heuristics collapse under multi-round execution and building a closed-loop controller to guarantee utility preservation.

\subsection{Catastrophic Forgetting and Subspace Methods}
Sequential unlearning is structurally related to continual learning in reverse. While continual learning seeks to acquire new tasks without catastrophic forgetting of previous tasks \cite{b6}, unlearning must intentionally excise target data while preserving retained knowledge. Saha et al. introduced Gradient Projection Memory (GPM) to project training steps orthogonal to past gradient subspaces \cite{b16}. SOM adapts the concept of orthonormal subspace tracking, but uses it as an observational sensor to drive closed-loop routing rather than as an optimization constraint.

\subsection{Auditing and Verification}
Membership Inference Attacks (MIA) \cite{b9}, \cite{b17} serve as standard empirical tools for evaluating residual private data leakage. We incorporate MIA AUC as one of three simultaneous gates in our closed-loop verifier, alongside validation forget accuracy and retain utility drop thresholds.

\section{Problem Formulation}
Let a base model $\theta_0 \in \mathbb{R}^d$ be trained on a dataset $D$ with $C$ classes. Over time, an ordered sequence of deletion requests $D_f^{(1)}, D_f^{(2)}, \dots, D_f^{(T)}$ arrives. At round $t$, the remaining retain dataset is:
\begin{equation}
D_r^{(t)} = D \setminus \bigcup_{\tau=1}^t D_f^{(\tau)}
\end{equation}
At each round, the controller selects an unlearning primitive $a_t \in \{1, 2, 3\}$, updating the model weights according to:
\begin{equation}
\theta_t = \mathcal{U}_{a_t}\left(\theta_{t-1}, D_f^{(t)}, D_r^{(t)}\right) \label{eq:update}
\end{equation}
A candidate model $\theta_t$ is acceptable if and only if three verification conditions hold simultaneously:
\begin{align}
\Acc\left(D_{f,\text{val}}^{(t)}; \theta_t\right) &\le \frac{1}{C}, \label{eq:verif1} \\
\Acc\left(D_{r,\text{anchor}}^{(t)}; \theta_{t-1}\right) - \Acc\left(D_{r,\text{anchor}}^{(t)}; \theta_t\right) &\le \varepsilon, \label{eq:verif2} \\
\left|\AUC_{\text{MIA}}(\theta_t) - 0.5\right| &\le \delta \label{eq:verif3}
\end{align}
where $1/C$ is the random chance threshold (10.0\% for CIFAR-10), $\varepsilon = 3.0\%$ is the maximum tolerable retain accuracy drop per round relative to the baseline checkpoint, and $\delta = 0.05$ bounds membership inference attack advantage. Exact retraining always satisfies \eqref{eq:verif1}--\eqref{eq:verif3} at maximum compute cost; the objective of ASUC-SOM is to minimize total execution time across the sequence while strictly enforcing verification at every round.

\section{The ASUC-SOM Framework}

\subsection{System Overview}
Fig.~\ref{fig:arch} illustrates the sequential unlearning pipeline. When batch $D_f^{(t)}$ arrives, the request profiler computes batch cardinality and class entropy. The Subspace Overlap Monitor (SOM) evaluates gradient projection overlap against historical updates. The tiered router dispatches the request to the optimal primitive, and the verification gate validates amnesia and utility before committing the checkpoint.

\begin{figure}[htbp]
\centering
\includegraphics[width=\columnwidth]{fig1_sequential_accuracy.png}
\caption{Processing path for sequential deletion requests in ASUC-SOM. A failed verification check rolls back model weights to $\theta_{t-1}$ and escalates the request tier.}
\label{fig:arch}
\end{figure}

\subsection{Request Profiling and Class Entropy}
The request profiler extracts two invariant metrics from $D_f^{(t)}$: batch cardinality ratio $\rho = |D_f^{(t)}| / |D_r^{(t)}|$ and normalized Shannon class entropy:
\begin{equation}
h = -\frac{1}{\log C} \sum_{c=1}^C p(c) \log p(c) \in [0, 1] \label{eq:entropy}
\end{equation}
where $p(c)$ is the empirical fraction of $D_f^{(t)}$ belonging to class $c$. When $h \approx 1.0$, deletion requests are uniformly distributed across classes. When $h < 0.40$, requests are class-concentrated, presenting high risk of destroying shared convolutional representations if unweighted dampening is used.

\subsection{Subspace Overlap Monitor (SOM)}
Scalar parameter distance $\|\theta_t - \theta_0\|_2$ measures total weight displacement but cannot determine whether updates disturb already-modified feature channels. SOM maintains an orthonormal basis $Q_k \in \mathbb{R}^{d \times k}$ spanning previous gradient directions. Given the incoming forget-loss gradient $g_t = \nabla_\theta \mathcal{L}(D_f^{(t)}; \theta_{t-1})$, the projection overlap score is defined as:
\begin{equation}
\text{SOM}(g_t, Q_k) = \frac{\|Q_k^T g_t\|_2}{\|g_t\|_2} \in [0, 1] \label{eq:som}
\end{equation}
Because SOM is computed from the gradient prior to modifying weights, it predicts parameter fatigue before damage occurs. After each successful round, the verified update $\Delta \theta = \theta_t - \theta_{t-1}$ is orthogonalized via Gram-Schmidt and appended to $Q_k$. When basis rank exceeds $k_{\max} = 50$, truncated SVD preserves the dominant orthogonal energy components.

\subsection{Multi-Tier Dynamic Router}
The router maps the tuple $(s, h)$ to execution tier $a \in \{1, 2, 3\}$, where $s = \text{SOM}(g_t, Q_k)$:
\begin{equation}
a = \begin{cases}
3, & \text{if } s \ge \tau_h \\
2, & \text{if } \tau_l \le s < \tau_h \text{ or } h < \tau_H \\
1, & \text{otherwise}
\end{cases} \label{eq:router}
\end{equation}
with default threshold parameters $\tau_l = 0.40$, $\tau_h = 0.80$, and $\tau_H = 0.40$. Table~\ref{tab:policy} outlines the routing policy and speedup factors.

\begin{table}[htbp]
\caption{ASUC-SOM Tiered Routing Policy}
\label{tab:policy}
\centering
\resizebox{\columnwidth}{!}{%
\begin{tabular}{cclc}
\toprule
\textbf{Tier} & \textbf{Trigger Condition} & \textbf{Unlearning Primitive} & \textbf{Speedup} \\
\midrule
1 & $s < 0.40 \land h \ge 0.40$ & Synaptic Dampening (SSD) \cite{b4} & $\approx 150\times$ \\
2 & $s \ge 0.40 \lor h < 0.40$ & Saliency Masking (SalUn) \cite{b5} & $\approx 115\times$ \\
3 & $s \ge 0.80$ / Escalation & Exact Retraining on $D_r^{(t)}$ & $1\times$ \\
\bottomrule
\end{tabular}%
}
\end{table}

\subsection{Closed-Loop Verification and Rollback Gate}
Following primitive execution, candidate model $\theta_t^*$ is evaluated on held-out validation sets. If all conditions in \eqref{eq:verif1}--\eqref{eq:verif3} pass, the model is committed. If any check fails, weights are restored to pre-round checkpoint $\theta_{t-1}$, the tier index is incremented ($a \leftarrow a + 1$), and the request is re-executed under regularized or exact retraining tiers.

\subsection{Tamper-Evident Compliance Audit Trail}
The compliance engine records a hash-chained JSON audit record for every batch:
\begin{equation}
\begin{aligned}
H_t = \text{SHA-256}\big( &H_{t-1} \,\|\, t \,\|\, |D_f^{(t)}| \,\|\, s_t \\
&\|\, a_t \,\|\, \Acc_f \,\|\, \Acc_r \,\|\, \text{Hash}(\theta_t)\big)
\end{aligned}
\end{equation}
providing verifiable proof of GDPR Article 17 amnesia compliance.

\begin{algorithm}[htbp]
\caption{ASUC-SOM Sequential Unlearning Routine}
\label{alg:asuc}
\small
\begin{algorithmic}[1]
\REQUIRE Model $\theta_{t-1}$, basis $Q_k$, forget set $D_f^{(t)}$, retain set $D_r^{(t)}$
\ENSURE Updated model $\theta_t$, updated basis $Q_k$, audit record $H_t$
\STATE $\theta_{\text{prev}} \leftarrow \theta_{t-1}$ \hfill \COMMENT{Save rollback checkpoint}
\STATE Compute cardinality $n = |D_f^{(t)}|$ and class entropy $h$ via \eqref{eq:entropy}
\STATE Compute gradient $g_t = \nabla_\theta \mathcal{L}(D_f^{(t)}; \theta_{\text{prev}})$
\STATE Compute subspace overlap $s = \|Q_k^T g_t\|_2 / \|g_t\|_2$ via \eqref{eq:som}
\STATE Determine primitive tier $a \leftarrow \text{Route}(s, h)$ via \eqref{eq:router}
\REPEAT
    \STATE $\theta^* \leftarrow \text{ExecutePrimitive}(a, \theta_{\text{prev}}, D_f^{(t)}, D_r^{(t)})$
    \IF{$\text{Verify}(\theta^*)$ passes via \eqref{eq:verif1}--\eqref{eq:verif3}}
        \STATE \textbf{break}
    \ENDIF
    \STATE $\theta^* \leftarrow \theta_{\text{prev}}; \quad a \leftarrow a + 1$ \hfill \COMMENT{Rollback \& escalate}
\UNTIL{$a > 3$}
\IF{$a = 3$}
    \STATE Reset basis $Q_k \leftarrow \emptyset$ \hfill \COMMENT{Clean retraining reset}
\ELSE
    \STATE $Q_k \leftarrow \text{UpdateBasis}(Q_k, \theta^* - \theta_{\text{prev}})$
\ENDIF
\STATE Append hash-chained audit record $H_t$
\STATE \textbf{return} $\theta_t = \theta^*, Q_k$
\end{algorithmic}
\end{algorithm}

\section{Experimental Evaluation Design}
\textbf{Model \& Dataset:} We evaluate on CIFAR-10 \cite{b11} (50,000 training images, 10,000 test images, 10 classes) using a ResNet-18 architecture (11,173,962 parameters). The base model achieves 93.44\% test accuracy and 99.40\% train retain accuracy after 200 epochs of SGD optimization ($\eta=0.1$, momentum 0.9, weight decay $5\times 10^{-4}$, cosine annealing).

\textbf{Sequential Workloads:}
\begin{itemize}
    \item \textbf{Workload A (Uniform Random):} 20 deletion rounds of 500 randomly sampled images per round across all 10 classes (10,000 total deleted images = 20\% of training data).
    \item \textbf{Workload B (Class-Sequential):} 10 deletion rounds of 500 samples per round concentrated entirely within a single class ($h = 0.0$, 5,000 total deleted samples).
    \item \textbf{Workload C (Dynamic High-Loss):} 10 deletion rounds of 500 samples prioritized by top cross-entropy loss values.
\end{itemize}

\textbf{Baselines:} We benchmark against Gradient Ascent (GA), Selective Synaptic Dampening (SSD) \cite{b4}, Saliency Unlearning (SalUn) \cite{b5}, and Exact Retraining from scratch on $D_r^{(t)}$.

\section{Results and Discussion}

\subsection{Cost Arithmetic}
Relative to exact retraining, SSD (Tier 1) delivers $\approx 150\times$ speedup and SalUn (Tier 2) delivers $\approx 115\times$ speedup. Tier 3 costs $1\times$ by definition. If $f_1, f_2, f_3$ denote the empirical fractions of rounds executed at each tier, the average cost per round $\bar{c}$ and speedup are:
\begin{equation}
\bar{c} = \frac{f_1}{150} + \frac{f_2}{115} + f_3, \quad \text{Speedup} = \frac{1}{\bar{c}} \label{eq:cost}
\end{equation}
If only one round in twenty escalates to exact retraining ($f_3 = 0.05$), the stream-level speedup is $\approx 17.5\times$. Thus, keeping exact retraining rare while preventing rank collapse is the central operational objective of the SOM monitor.

\subsection{Sequential Performance over 20 Rounds (Workload A)}
Table~\ref{tab:results} summarizes comparative performance across 20 sequential deletion rounds under Workload A.

\begin{table*}[t]
\caption{Sequential Machine Unlearning Benchmark Matrix on CIFAR-10 ResNet-18 (Workload A, 20 Rounds, 10,000 Deleted Samples)}
\label{tab:results}
\centering
\resizebox{\textwidth}{!}{%
\begin{tabular}{lcccccc}
\toprule
\textbf{Method} & \textbf{Collapse Round} & \textbf{Final Retain Acc. (\%)} & \textbf{Final Test Acc. (\%)} & \textbf{Forget Acc. (\%)} & \textbf{MIA AUC} & \textbf{Total Time (s)} \\
\midrule
Gradient Ascent (GA) & None (Failed to forget) & 99.47 & 93.19 & 99.80 & 0.998 & 70.7 \\
SSD (Foster et al., 2024) \cite{b4} & \textbf{Round 3 (Collapse)} & 9.96 & 10.00 & 10.00 & 0.500 & 73.2 \\
SalUn (Fan et al., 2024) \cite{b5} & None ($>20$ Stable) & 95.03 & 88.03 & 10.40 & 0.538 & 303.5 \\
Exact Retraining (Gold Standard) & None ($>20$ Stable) & 99.10 & 92.70 & 9.20 & 0.512 & $\approx 21,100.0$ \\
\textbf{ASUC-SOM (Ours)} & \textbf{None ($>20$ Stable)} & \textbf{97.19} & \textbf{90.89} & \textbf{10.20} & \textbf{0.546} & \textbf{15,655.4} \\
\bottomrule
\end{tabular}%
}
\end{table*}

\begin{figure}[htbp]
\centering
\includegraphics[width=\columnwidth]{fig2_som_dynamics_drift.png}
\caption{Subspace Overlap Metric (SOM) trajectory and cumulative parameter drift $\|\Delta\theta\|_2$ across sequential deletion rounds. Green, Amber, and Red health states govern dynamic tier dispatch.}
\label{fig:som_drift}
\end{figure}

As shown in Table~\ref{tab:results}, static SSD undergoes catastrophic representational collapse at Round 3. Retain accuracy drops from 99.53\% at Round 1 to 49.42\% at Round 2, and collapses to 9.96\% (random guessing) by Round 3. This occurs because repeated Fisher dampening indiscriminately suppresses overlapping parameter dimensions without regularizing against retain utility loss. Gradient Ascent fails to forget target data, exhibiting 99.80\% residual forget accuracy.

In contrast, ASUC-SOM continuously tracks directional gradient overlap via SOM (Fig.~\ref{fig:som_drift}). When the overlap score enters the Amber caution zone ($s \ge 0.40$), the controller dynamically shifts execution to Tier 2 (SalUn), successfully preserving 97.19\% retain accuracy and 90.89\% test generalization across all 20 rounds.

\begin{table}[htbp]
\caption{Architectural Comparison: Static Unlearning vs. ASUC-SOM}
\label{tab:comparison}
\centering
\resizebox{\columnwidth}{!}{%
\begin{tabular}{lll}
\toprule
\textbf{Dimension} & \textbf{Static Baselines (SSD / SalUn / GA)} & \textbf{ASUC-SOM (Ours)} \\
\midrule
Operational Scope & Single isolated request & Continuous stream (20+ rounds) \\
Routing Strategy & Single fixed heuristic & Dynamic three-tier routing \\
Geometry Tracking & Not tracked & Orthonormal basis $Q_k$, SOM score \\
Verification & Open-loop (no checks) & Closed-loop verify, rollback, escalate \\
Auditability & Ad-hoc benchmark logs & Hash-chained per-batch audit records \\
\bottomrule
\end{tabular}%
}
\end{table}

\subsection{Multi-Workload Stress Testing (Workloads B and C)}
Under Workload B (Class-Sequential, 10 rounds), zero class entropy ($h = 0.0$) is immediately flagged by the profiler. The controller bypasses single-pass SSD in favor of Tier 2 (SalUn), sustaining 98.26\% retain accuracy and 92.33\% test accuracy. Under Workload C (Dynamic High-Loss), ASUC-SOM sustains 99.11\% retain accuracy and 90.75\% test accuracy.

\section{Limitations}
We identify four key operational boundaries:
\begin{enumerate}
    \item \textbf{Basis Dimensionality Scaling:} Maintaining basis $Q_k$ scales as $\mathcal{O}(d \cdot k)$ in memory. For billion-parameter vision-language models, low-rank parameter adapters (LoRA) or randomized sketch projections will be necessary.
    \item \textbf{Anchor Set Isolation:} The closed-loop verification gate requires a curated retain anchor set $D_{r,\text{anchor}}$ that must remain strictly isolated from deletion candidates.
    \item \textbf{Empirical Thresholds:} Thresholds $\tau_l=0.40$ and $\tau_h=0.80$ were calibrated on vision architectures and may require adaptive learning rate tuning for transformers.
    \item \textbf{Differential Privacy Bounds:} While empirical membership inference AUC is bounded at $\le 0.55$, formal $(\epsilon, \delta)$-differential privacy guarantees across sequential rounds remain an open theoretical avenue.
\end{enumerate}

\section{Conclusion}
Regulatory machine unlearning is fundamentally an ongoing, sequential process. Static one-pass heuristics fail under repeated execution due to parameter fatigue and subspace overlap collisions. ASUC-SOM resolves this challenge by formulating sequential unlearning as a closed-loop geometric control problem. By monitoring subspace overlap via SOM, dynamically dispatching deletion requests across adaptive primitives, and enforcing verifiable amnesia through automated rollback gates, ASUC-SOM eliminates catastrophic rank collapse across 20 sequential deletion rounds on CIFAR-10 ResNet-18 while achieving up to 148$\times$ compute speedup over full retraining.

\begin{thebibliography}{00}
\bibitem{b1} European Parliament and Council, ``Regulation (EU) 2016/679 (General Data Protection Regulation),'' \textit{Official Journal of the European Union}, Art. 17, 2016.
\bibitem{b2} Y. Cao and J. Yang, ``Towards making systems forget with machine unlearning,'' in \textit{Proc. IEEE Symp. Security and Privacy}, 2015, pp. 463--480.
\bibitem{b3} L. Bourtoule et al., ``Machine unlearning,'' in \textit{Proc. IEEE Symp. Security and Privacy}, 2021, pp. 141--159.
\bibitem{b4} J. Foster, S. Schoepf, and A. Brintrup, ``Fast machine unlearning without retraining through selective synaptic dampening,'' in \textit{Proc. AAAI Conf. Artificial Intelligence}, vol. 38, 2024, pp. 12043--12051.
\bibitem{b5} C. Fan, J. Liu, Y. Zhang, E. Wong, D. Wei, and S. Liu, ``SalUn: Empowering machine unlearning via gradient-based weight saliency in both image classification and generation,'' in \textit{Proc. Int. Conf. Learning Representations}, 2024.
\bibitem{b6} J. Kirkpatrick et al., ``Overcoming catastrophic forgetting in neural networks,'' \textit{Proc. Natl. Acad. Sci. USA}, vol. 114, no. 13, pp. 3521--3526, 2017.
\bibitem{b7} A. Golatkar, A. Achille, and S. Soatto, ``Eternal sunshine of the spotless net: Selective forgetting in deep networks,'' in \textit{Proc. IEEE/CVF Conf. Computer Vision and Pattern Recognition}, 2020, pp. 9304--9312.
\bibitem{b8} A. Thudi, G. Deza, V. Chandrasekaran, and N. Papernot, ``Unrolling SGD: Understanding factors influencing machine unlearning,'' in \textit{Proc. IEEE Eur. Symp. Security and Privacy}, 2022, pp. 303--319.
\bibitem{b9} R. Shokri, M. Stronati, C. Song, and V. Shmatikov, ``Membership inference attacks against machine learning models,'' in \textit{Proc. IEEE Symp. Security and Privacy}, 2017, pp. 3--18.
\bibitem{b10} V. S. Chundawat, A. K. Tarun, M. Mandal, and M. Kankanhalli, ``Can bad teaching induce forgetting? Unlearning in deep networks using an incompetent teacher,'' in \textit{Proc. AAAI Conf. Artificial Intelligence}, vol. 37, 2023.
\bibitem{b11} A. Krizhevsky, ``Learning multiple layers of features from tiny images,'' Univ. Toronto, Tech. Rep., 2009.
\bibitem{b12} V. Gupta, C. Jung, S. Neel, A. Roth, S. Sharifi-Malvajerdi, and C. Waites, ``Adaptive machine unlearning,'' in \textit{Proc. Adv. Neural Information Processing Systems}, vol. 34, 2021.
\bibitem{b13} L. Graves, V. Nagisetty, and V. Ganesh, ``Amnesiac machine learning,'' in \textit{Proc. AAAI Conf. Artificial Intelligence}, vol. 35, 2021, pp. 11516--11524.
\bibitem{b14} A. Ginart, M. Guan, G. Valiant, and J. Zou, ``Making AI forget you: Data deletion in machine learning,'' in \textit{Proc. Adv. Neural Information Processing Systems}, vol. 32, 2019.
\bibitem{b15} A. Sekhari, J. Acharya, G. Kamath, and A. T. Suresh, ``Remember what you want to forget: Algorithms for machine unlearning,'' in \textit{Proc. Adv. Neural Information Processing Systems}, vol. 34, 2021.
\bibitem{b16} G. Saha, I. Garg, and K. Roy, ``Gradient projection memory for continual learning,'' in \textit{Proc. Int. Conf. Learning Representations}, 2021.
\bibitem{b17} N. Carlini, S. Chien, M. Nasr, S. Song, A. Terzis, and F. Tramèr, ``Membership inference attacks from first principles,'' in \textit{Proc. IEEE Symp. Security and Privacy}, 2022, pp. 1897--1914.
\end{thebibliography}

\end{document}
